%%%%%%%%%%%%%%%%%%%%%%%%%%%%%%%%%%%%%%%%%%%%%%%%%%%%%%%%%%%%
% 10.7 DYNAMIC PROGRAMMING
% The Composition of Advanced Mathematics
%%%%%%%%%%%%%%%%%%%%%%%%%%%%%%%%%%%%%%%%%%%%%%%%%%%%%%%%%%%%

\section{Dynamic Programming}
\label{sec:dynamic-programming}

\prerequisite{
Optimization fundamentals, combinatorial optimization, recursion,
sequences, discrete mathematics, probability, graph theory, basic
algorithms, and elementary stochastic processes.
}

\learningobjectives{
By the end of this section, the reader should be able to:
\begin{itemize}
    \item define dynamic programming;
    \item identify optimal substructure;
    \item identify overlapping subproblems;
    \item distinguish recursive search from dynamic programming;
    \item define states, actions, stages, and transitions;
    \item formulate deterministic finite-horizon dynamic programs;
    \item derive Bellman recursions;
    \item explain the principle of optimality;
    \item distinguish top-down memoization and bottom-up tabulation;
    \item understand state-space design;
    \item recognize dimensionality of state spaces;
    \item formulate shortest-path problems as dynamic programs;
    \item solve the \(0\)-\(1\) knapsack problem by dynamic programming;
    \item solve resource-allocation problems by dynamic programming;
    \item formulate sequence and alignment problems;
    \item understand longest common subsequence recurrences;
    \item understand edit-distance recurrences;
    \item recognize interval dynamic programming;
    \item understand matrix-chain multiplication;
    \item recognize subset dynamic programming;
    \item formulate the Held--Karp algorithm for the traveling salesperson
          problem;
    \item understand pseudo-polynomial complexity;
    \item distinguish polynomial and pseudo-polynomial algorithms;
    \item understand infinite-horizon dynamic programming conceptually;
    \item formulate deterministic Bellman equations;
    \item formulate stochastic Bellman equations;
    \item understand expected-value recursions;
    \item recognize Markov decision processes;
    \item distinguish state-value and action-value functions;
    \item understand discounted dynamic programming;
    \item understand finite-horizon stochastic control;
    \item understand policy evaluation and policy improvement conceptually;
    \item recognize value iteration conceptually;
    \item understand the curse of dimensionality;
    \item recognize state aggregation and approximation methods;
    \item understand approximate dynamic programming conceptually;
    \item connect dynamic programming with optimal control,
          reinforcement learning, shortest paths, scheduling, inventory,
          and stochastic optimization.
\end{itemize}
}

%%%%%%%%%%%%%%%%%%%%%%%%%%%%%%%%%%%%%%%%%%%%%%%%%%%%%%%%%%%%
\subsection{What Is Dynamic Programming?}
\label{subsec:what-is-dynamic-programming}
%%%%%%%%%%%%%%%%%%%%%%%%%%%%%%%%%%%%%%%%%%%%%%%%%%%%%%%%%%%%

Dynamic programming is an optimization framework for problems that can be
decomposed into smaller, related subproblems.

Rather than solving every possible sequence of decisions independently,
dynamic programming stores and reuses the values of recurring
subproblems.

A general idea is

\[
\boxed{
\text{large optimization problem}
\rightarrow
\text{smaller optimization problems}
\rightarrow
\text{Bellman recursion}.
}
\]

%%%%%%%%%%%%%%%%%%%%%%%%%%%%%%%%%%%%%%%%%%%%%%%%%%%%%%%%%%%%
\subsection{Two Fundamental Properties}
\label{subsec:dynamic-programming-properties}
%%%%%%%%%%%%%%%%%%%%%%%%%%%%%%%%%%%%%%%%%%%%%%%%%%%%%%%%%%%%

Dynamic programming is especially useful when a problem has:

\[
\boxed{
\text{optimal substructure},
}
\]

and

\[
\boxed{
\text{overlapping subproblems}.
}
\]

Optimal substructure means an optimal solution can be constructed from
optimal solutions to smaller problems.

Overlapping subproblems means the same smaller problems appear repeatedly.

%%%%%%%%%%%%%%%%%%%%%%%%%%%%%%%%%%%%%%%%%%%%%%%%%%%%%%%%%%%%
\subsection{Optimal Substructure}
\label{subsec:optimal-substructure}
%%%%%%%%%%%%%%%%%%%%%%%%%%%%%%%%%%%%%%%%%%%%%%%%%%%%%%%%%%%%

Suppose the best solution from state \(s\) begins with action \(a\).

After taking that action, the remaining decisions must form an optimal
solution for the resulting state.

Otherwise the original solution could be improved.

This principle is the foundation of Bellman's principle of optimality.

%%%%%%%%%%%%%%%%%%%%%%%%%%%%%%%%%%%%%%%%%%%%%%%%%%%%%%%%%%%%
\subsection{Overlapping Subproblems}
\label{subsec:overlapping-subproblems}
%%%%%%%%%%%%%%%%%%%%%%%%%%%%%%%%%%%%%%%%%%%%%%%%%%%%%%%%%%%%

A naive recursion may solve the same subproblem many times.

For example, Fibonacci recursion uses

\[
F_n
=
F_{n-1}
+
F_{n-2}.
\]

Computing \(F_{n-1}\) and \(F_{n-2}\) recursively causes many repeated
evaluations.

Dynamic programming stores already-computed values.

%%%%%%%%%%%%%%%%%%%%%%%%%%%%%%%%%%%%%%%%%%%%%%%%%%%%%%%%%%%%
\subsection{Memoization}
\label{subsec:dynamic-programming-memoization}
%%%%%%%%%%%%%%%%%%%%%%%%%%%%%%%%%%%%%%%%%%%%%%%%%%%%%%%%%%%%

Memoization is the top-down form of dynamic programming.

One writes the recursive problem naturally and stores the value of each
state after it is first computed.

If that state appears again, the stored value is reused.

Thus,

\[
\boxed{
\text{recursion}
+
\text{memory}
=
\text{memoization}.
}
\]

%%%%%%%%%%%%%%%%%%%%%%%%%%%%%%%%%%%%%%%%%%%%%%%%%%%%%%%%%%%%
\subsection{Tabulation}
\label{subsec:dynamic-programming-tabulation}
%%%%%%%%%%%%%%%%%%%%%%%%%%%%%%%%%%%%%%%%%%%%%%%%%%%%%%%%%%%%

Tabulation is the bottom-up form of dynamic programming.

One identifies an order in which smaller states can be solved before larger
states.

The value table is filled iteratively.

Thus,

\[
\boxed{
\text{small states}
\rightarrow
\text{larger states}
\rightarrow
\text{final solution}.
}
\]

%%%%%%%%%%%%%%%%%%%%%%%%%%%%%%%%%%%%%%%%%%%%%%%%%%%%%%%%%%%%
\subsection{Memoization Versus Tabulation}
\label{subsec:memoization-versus-tabulation}
%%%%%%%%%%%%%%%%%%%%%%%%%%%%%%%%%%%%%%%%%%%%%%%%%%%%%%%%%%%%

Memoization:

\[
\boxed{
\text{top-down},
}
\]

and computes only states actually reached.

Tabulation:

\[
\boxed{
\text{bottom-up},
}
\]

and often computes states in a regular predetermined order.

Both implement the same underlying recurrence when formulated correctly.

%%%%%%%%%%%%%%%%%%%%%%%%%%%%%%%%%%%%%%%%%%%%%%%%%%%%%%%%%%%%
\subsection{State}
\label{subsec:dynamic-programming-state}
%%%%%%%%%%%%%%%%%%%%%%%%%%%%%%%%%%%%%%%%%%%%%%%%%%%%%%%%%%%%

A state contains all information from the past that is needed to make
optimal future decisions.

A state should summarize the relevant history without storing unnecessary
detail.

Examples include:

\[
\boxed{
\text{current location},
}
\]

\[
\boxed{
\text{remaining capacity},
}
\]

\[
\boxed{
\text{time period},
}
\]

or

\[
\boxed{
\text{current inventory}.
}
\]

%%%%%%%%%%%%%%%%%%%%%%%%%%%%%%%%%%%%%%%%%%%%%%%%%%%%%%%%%%%%
\subsection{Actions}
\label{subsec:dynamic-programming-actions}
%%%%%%%%%%%%%%%%%%%%%%%%%%%%%%%%%%%%%%%%%%%%%%%%%%%%%%%%%%%%

At state \(s\), the decision maker chooses an action

\[
\boxed{
a\in A(s),
}
\]

where \(A(s)\) is the feasible action set.

The action determines:

\[
\boxed{
\text{immediate cost or reward},
}
\]

and

\[
\boxed{
\text{next state}.
}
\]

%%%%%%%%%%%%%%%%%%%%%%%%%%%%%%%%%%%%%%%%%%%%%%%%%%%%%%%%%%%%
\subsection{Transition Function}
\label{subsec:dynamic-programming-transition}
%%%%%%%%%%%%%%%%%%%%%%%%%%%%%%%%%%%%%%%%%%%%%%%%%%%%%%%%%%%%

In a deterministic problem, the next state can be written

\[
\boxed{
s'
=
T(s,a).
}
\]

Thus a dynamic decision has the form

\[
\boxed{
s
\xrightarrow{a}
T(s,a).
}
\]

In stochastic problems, the next state is random.

%%%%%%%%%%%%%%%%%%%%%%%%%%%%%%%%%%%%%%%%%%%%%%%%%%%%%%%%%%%%
\subsection{Stage}
\label{subsec:dynamic-programming-stage}
%%%%%%%%%%%%%%%%%%%%%%%%%%%%%%%%%%%%%%%%%%%%%%%%%%%%%%%%%%%%

Many dynamic programs evolve through stages

\[
\boxed{
t=0,1,\ldots,T.
}
\]

The stage may represent:

\[
\boxed{
\text{time},
}
\]

\[
\boxed{
\text{position in a sequence},
}
\]

\[
\boxed{
\text{number of items processed},
}
\]

or

\[
\boxed{
\text{decision depth}.
}
\]

%%%%%%%%%%%%%%%%%%%%%%%%%%%%%%%%%%%%%%%%%%%%%%%%%%%%%%%%%%%%
\subsection{Value Function}
\label{subsec:dynamic-programming-value-function}
%%%%%%%%%%%%%%%%%%%%%%%%%%%%%%%%%%%%%%%%%%%%%%%%%%%%%%%%%%%%

The value function records the best achievable future objective from a
given state.

For minimization,

\[
\boxed{
V_t(s)
=
\text{minimum future cost from state }s\text{ at stage }t.
}
\]

For maximization,

\[
\boxed{
V_t(s)
=
\text{maximum future reward from state }s.
}
\]

%%%%%%%%%%%%%%%%%%%%%%%%%%%%%%%%%%%%%%%%%%%%%%%%%%%%%%%%%%%%
\subsection{Bellman's Principle of Optimality}
\label{subsec:bellman-principle-optimality}
%%%%%%%%%%%%%%%%%%%%%%%%%%%%%%%%%%%%%%%%%%%%%%%%%%%%%%%%%%%%

\begin{definitionbox}{Bellman's Principle of Optimality}
An optimal policy has the property that, regardless of the initial state
and first decision, the remaining decisions must constitute an optimal
policy for the state resulting from that first decision.
\end{definitionbox}

This principle permits a global optimization problem to be written
recursively.

%%%%%%%%%%%%%%%%%%%%%%%%%%%%%%%%%%%%%%%%%%%%%%%%%%%%%%%%%%%%
\subsection{Finite-Horizon Bellman Recursion}
\label{subsec:finite-horizon-bellman-recursion}
%%%%%%%%%%%%%%%%%%%%%%%%%%%%%%%%%%%%%%%%%%%%%%%%%%%%%%%%%%%%

Suppose:

\[
c_t(s,a)
=
\text{cost at stage }t,
\]

and

\[
s'
=
T_t(s,a).
\]

Then a deterministic finite-horizon recursion is

\[
\boxed{
V_t(s)
=
\min_{a\in A_t(s)}
\left[
c_t(s,a)
+
V_{t+1}
\left(
T_t(s,a)
\right)
\right].
}
\]

A terminal cost may be

\[
\boxed{
V_T(s)=g(s).
}
\]

%%%%%%%%%%%%%%%%%%%%%%%%%%%%%%%%%%%%%%%%%%%%%%%%%%%%%%%%%%%%
\subsection{Backward Induction}
\label{subsec:backward-induction}
%%%%%%%%%%%%%%%%%%%%%%%%%%%%%%%%%%%%%%%%%%%%%%%%%%%%%%%%%%%%

Finite-horizon dynamic programming is usually solved backward.

First compute

\[
V_T.
\]

Then compute

\[
V_{T-1},
\]

followed by

\[
V_{T-2},
\]

until reaching

\[
V_0.
\]

This is called backward induction.

%%%%%%%%%%%%%%%%%%%%%%%%%%%%%%%%%%%%%%%%%%%%%%%%%%%%%%%%%%%%
\subsection{Recovering the Optimal Policy}
\label{subsec:recovering-optimal-policy}
%%%%%%%%%%%%%%%%%%%%%%%%%%%%%%%%%%%%%%%%%%%%%%%%%%%%%%%%%%%%

The Bellman recursion gives the optimal objective value.

To recover the decisions, store an optimizer:

\[
\boxed{
\pi_t^*(s)
\in
\operatorname*{arg\,min}_{a\in A_t(s)}
\left[
c_t(s,a)
+
V_{t+1}
(
T_t(s,a)
)
\right].
}
\]

The mapping

\[
\pi_t^*
\]

is an optimal policy.

%%%%%%%%%%%%%%%%%%%%%%%%%%%%%%%%%%%%%%%%%%%%%%%%%%%%%%%%%%%%
\subsection{Worked Example: Two-Stage Decision}
\label{subsec:worked-two-stage-dp}
%%%%%%%%%%%%%%%%%%%%%%%%%%%%%%%%%%%%%%%%%%%%%%%%%%%%%%%%%%%%

\begin{workedexamplebox}{Simple Backward Induction}

Suppose at the final stage the costs are

\[
V_2(A)=3,
\qquad
V_2(B)=7.
\]

At stage \(1\), from state \(S\), there are two actions.

Action \(a_1\) costs \(4\) and moves to \(A\).

Action \(a_2\) costs \(1\) and moves to \(B\).

Then

\[
V_1(S)
=
\min
\{
4+3,
1+7
\}.
\]

Thus,

\[
V_1(S)
=
\min
\{
7,8
\}.
\]

Therefore,

\begin{answerbox}
\[
\boxed{
V_1(S)=7,
}
\]

and the optimal action is

\[
\boxed{
a_1.
}
\]
\end{answerbox}

\end{workedexamplebox}

%%%%%%%%%%%%%%%%%%%%%%%%%%%%%%%%%%%%%%%%%%%%%%%%%%%%%%%%%%%%
\subsection{State-Space Graph Interpretation}
\label{subsec:state-space-graph}
%%%%%%%%%%%%%%%%%%%%%%%%%%%%%%%%%%%%%%%%%%%%%%%%%%%%%%%%%%%%

A dynamic program can often be represented as a directed graph.

Each vertex represents a state.

Each directed edge represents a feasible transition.

An edge may have cost

\[
\boxed{
c(s,a).
}
\]

Then dynamic programming finds an optimal path through the state-space
graph.

%%%%%%%%%%%%%%%%%%%%%%%%%%%%%%%%%%%%%%%%%%%%%%%%%%%%%%%%%%%%
\subsection{Acyclic State Graphs}
\label{subsec:acyclic-state-graphs}
%%%%%%%%%%%%%%%%%%%%%%%%%%%%%%%%%%%%%%%%%%%%%%%%%%%%%%%%%%%%

If the state-space graph is acyclic, states can be processed in
topological order.

The value recurrence becomes

\[
\boxed{
V(v)
=
\min_{(v,w)\in E}
\left[
c_{vw}+V(w)
\right].
}
\]

This is a shortest-path dynamic program on a directed acyclic graph.

%%%%%%%%%%%%%%%%%%%%%%%%%%%%%%%%%%%%%%%%%%%%%%%%%%%%%%%%%%%%
\subsection{Shortest Path as Dynamic Programming}
\label{subsec:shortest-path-dynamic-programming}
%%%%%%%%%%%%%%%%%%%%%%%%%%%%%%%%%%%%%%%%%%%%%%%%%%%%%%%%%%%%

Let

\[
V(v)
\]

denote the shortest remaining distance from vertex \(v\) to destination
\(t\).

Then

\[
\boxed{
V(t)=0,
}
\]

and

\[
\boxed{
V(v)
=
\min_{(v,w)\in E}
\left[
c_{vw}+V(w)
\right].
}
\]

This is precisely a Bellman recursion.

%%%%%%%%%%%%%%%%%%%%%%%%%%%%%%%%%%%%%%%%%%%%%%%%%%%%%%%%%%%%
\subsection{Bellman--Ford Connection}
\label{subsec:bellman-ford-dynamic-programming}
%%%%%%%%%%%%%%%%%%%%%%%%%%%%%%%%%%%%%%%%%%%%%%%%%%%%%%%%%%%%

Let

\[
V_k(v)
\]

denote the minimum cost of a path from source \(s\) to \(v\) using at most
\(k\) edges.

Then

\[
\boxed{
V_k(v)
=
\min
\left\{
V_{k-1}(v),
\;
\min_{(u,v)\in E}
[
V_{k-1}(u)+c_{uv}
]
\right\}.
}
\]

This recurrence underlies the Bellman--Ford shortest-path algorithm.

%%%%%%%%%%%%%%%%%%%%%%%%%%%%%%%%%%%%%%%%%%%%%%%%%%%%%%%%%%%%
\subsection{Resource-Allocation Dynamic Programming}
\label{subsec:resource-allocation-dp}
%%%%%%%%%%%%%%%%%%%%%%%%%%%%%%%%%%%%%%%%%%%%%%%%%%%%%%%%%%%%

Suppose \(B\) units of a resource must be allocated among \(n\) activities.

Let

\[
x_i
\]

be the amount assigned to activity \(i\), and let

\[
r_i(x_i)
\]

be its reward.

We seek

\[
\boxed{
\max
\sum_{i=1}^{n}
r_i(x_i)
}
\]

subject to

\[
\boxed{
\sum_{i=1}^{n}
x_i
\leq
B.
}
\]

%%%%%%%%%%%%%%%%%%%%%%%%%%%%%%%%%%%%%%%%%%%%%%%%%%%%%%%%%%%%
\subsection{Resource-Allocation Recurrence}
\label{subsec:resource-allocation-recurrence}
%%%%%%%%%%%%%%%%%%%%%%%%%%%%%%%%%%%%%%%%%%%%%%%%%%%%%%%%%%%%

Let

\[
V_i(b)
\]

be the best reward obtained from the first \(i\) activities using at most
\(b\) units.

Then

\[
\boxed{
V_i(b)
=
\max_{0\leq x_i\leq b}
\left[
r_i(x_i)
+
V_{i-1}(b-x_i)
\right].
}
\]

The boundary condition is

\[
\boxed{
V_0(b)=0.
}
\]

%%%%%%%%%%%%%%%%%%%%%%%%%%%%%%%%%%%%%%%%%%%%%%%%%%%%%%%%%%%%
\subsection{0--1 Knapsack Dynamic Programming}
\label{subsec:knapsack-dynamic-programming}
%%%%%%%%%%%%%%%%%%%%%%%%%%%%%%%%%%%%%%%%%%%%%%%%%%%%%%%%%%%%

Suppose item \(i\) has value \(v_i\) and integer weight \(w_i\).

Let

\[
V(i,b)
\]

be the maximum value obtainable using the first \(i\) items with capacity
\(b\).

If

\[
w_i>b,
\]

item \(i\) cannot be selected:

\[
\boxed{
V(i,b)
=
V(i-1,b).
}
\]

If

\[
w_i\leq b,
\]

then

\[
\boxed{
V(i,b)
=
\max
\left\{
V(i-1,b),
\;
v_i
+
V(i-1,b-w_i)
\right\}.
}
\]

%%%%%%%%%%%%%%%%%%%%%%%%%%%%%%%%%%%%%%%%%%%%%%%%%%%%%%%%%%%%
\subsection{Knapsack Boundary Conditions}
\label{subsec:knapsack-boundary-conditions}
%%%%%%%%%%%%%%%%%%%%%%%%%%%%%%%%%%%%%%%%%%%%%%%%%%%%%%%%%%%%

The base cases are

\[
\boxed{
V(0,b)=0
}
\]

for every \(b\), and

\[
\boxed{
V(i,0)=0.
}
\]

The final optimal value is

\[
\boxed{
V(n,W).
}
\]

%%%%%%%%%%%%%%%%%%%%%%%%%%%%%%%%%%%%%%%%%%%%%%%%%%%%%%%%%%%%
\subsection{Worked Example: 0--1 Knapsack}
\label{subsec:worked-knapsack-dp}
%%%%%%%%%%%%%%%%%%%%%%%%%%%%%%%%%%%%%%%%%%%%%%%%%%%%%%%%%%%%

\begin{workedexamplebox}{Knapsack Recurrence}

Suppose three items have

\[
(w_1,v_1)=(2,4),
\]

\[
(w_2,v_2)=(3,5),
\]

\[
(w_3,v_3)=(4,7),
\]

with capacity

\[
W=5.
\]

For the first two items,

\[
V(2,5)
=
\max
\{
V(1,5),
5+V(1,2)
\}.
\]

Since

\[
V(1,5)=4,
\]

and

\[
V(1,2)=4,
\]

we obtain

\[
V(2,5)=9.
\]

For item \(3\),

\[
V(3,5)
=
\max
\{
V(2,5),
7+V(2,1)
\}.
\]

Since

\[
V(2,1)=0,
\]

\[
V(3,5)
=
\max
\{
9,7
\}
=
9.
\]

Therefore,

\begin{answerbox}
\[
\boxed{
V(3,5)=9.
}
\]

The optimal choice uses items \(1\) and \(2\).
\end{answerbox}

\end{workedexamplebox}

%%%%%%%%%%%%%%%%%%%%%%%%%%%%%%%%%%%%%%%%%%%%%%%%%%%%%%%%%%%%
\subsection{Knapsack Complexity}
\label{subsec:knapsack-dp-complexity}
%%%%%%%%%%%%%%%%%%%%%%%%%%%%%%%%%%%%%%%%%%%%%%%%%%%%%%%%%%%%

The table contains approximately

\[
nW
\]

states.

Each state requires constant work.

Thus the time complexity is

\[
\boxed{
O(nW).
}
\]

The memory complexity is also

\[
\boxed{
O(nW)
}
\]

for the full table.

%%%%%%%%%%%%%%%%%%%%%%%%%%%%%%%%%%%%%%%%%%%%%%%%%%%%%%%%%%%%
\subsection{Pseudo-Polynomial Complexity}
\label{subsec:dynamic-programming-pseudo-polynomial}
%%%%%%%%%%%%%%%%%%%%%%%%%%%%%%%%%%%%%%%%%%%%%%%%%%%%%%%%%%%%

Although

\[
O(nW)
\]

looks polynomial, the input contains \(W\) in binary.

The number of bits required to represent \(W\) is approximately

\[
\log_2W.
\]

Thus the running time is not polynomial in the formal input length.

It is called

\[
\boxed{
\text{pseudo-polynomial}.
}
\]

%%%%%%%%%%%%%%%%%%%%%%%%%%%%%%%%%%%%%%%%%%%%%%%%%%%%%%%%%%%%
\subsection{Memory Reduction in Knapsack}
\label{subsec:knapsack-memory-reduction}
%%%%%%%%%%%%%%%%%%%%%%%%%%%%%%%%%%%%%%%%%%%%%%%%%%%%%%%%%%%%

The recurrence uses only the previous item row.

Therefore the table can be reduced from

\[
O(nW)
\]

memory to

\[
\boxed{
O(W).
}
\]

For \(0\)-\(1\) knapsack, capacities should be updated in decreasing order
when using one array so that an item is not accidentally used multiple
times.

%%%%%%%%%%%%%%%%%%%%%%%%%%%%%%%%%%%%%%%%%%%%%%%%%%%%%%%%%%%%
\subsection{Unbounded Knapsack}
\label{subsec:unbounded-knapsack-dp}
%%%%%%%%%%%%%%%%%%%%%%%%%%%%%%%%%%%%%%%%%%%%%%%%%%%%%%%%%%%%

If each item can be selected repeatedly, the recurrence changes.

One form is

\[
\boxed{
V(b)
=
\max_{i:w_i\leq b}
\left[
v_i+V(b-w_i)
\right].
}
\]

This illustrates an important principle:

\[
\boxed{
\text{state recurrence depends on whether decisions can be reused}.
}
\]

%%%%%%%%%%%%%%%%%%%%%%%%%%%%%%%%%%%%%%%%%%%%%%%%%%%%%%%%%%%%
\subsection{Longest Common Subsequence}
\label{subsec:longest-common-subsequence}
%%%%%%%%%%%%%%%%%%%%%%%%%%%%%%%%%%%%%%%%%%%%%%%%%%%%%%%%%%%%

Given sequences

\[
X=x_1x_2\cdots x_m
\]

and

\[
Y=y_1y_2\cdots y_n,
\]

a common subsequence appears in both sequences while preserving order.

Let

\[
L(i,j)
\]

be the length of the longest common subsequence of the prefixes

\[
x_1\cdots x_i
\]

and

\[
y_1\cdots y_j.
\]

%%%%%%%%%%%%%%%%%%%%%%%%%%%%%%%%%%%%%%%%%%%%%%%%%%%%%%%%%%%%
\subsection{LCS Recurrence}
\label{subsec:lcs-recurrence}
%%%%%%%%%%%%%%%%%%%%%%%%%%%%%%%%%%%%%%%%%%%%%%%%%%%%%%%%%%%%

If

\[
x_i=y_j,
\]

then

\[
\boxed{
L(i,j)
=
1+L(i-1,j-1).
}
\]

If

\[
x_i\neq y_j,
\]

then

\[
\boxed{
L(i,j)
=
\max
\{
L(i-1,j),
L(i,j-1)
\}.
}
\]

Boundary conditions are

\[
\boxed{
L(i,0)=L(0,j)=0.
}
\]

%%%%%%%%%%%%%%%%%%%%%%%%%%%%%%%%%%%%%%%%%%%%%%%%%%%%%%%%%%%%
\subsection{LCS Complexity}
\label{subsec:lcs-complexity}
%%%%%%%%%%%%%%%%%%%%%%%%%%%%%%%%%%%%%%%%%%%%%%%%%%%%%%%%%%%%

There are

\[
(m+1)(n+1)
\]

states.

Each state requires constant work.

Thus the basic dynamic program runs in

\[
\boxed{
O(mn)
}
\]

time.

%%%%%%%%%%%%%%%%%%%%%%%%%%%%%%%%%%%%%%%%%%%%%%%%%%%%%%%%%%%%
\subsection{Edit Distance}
\label{subsec:edit-distance-dynamic-programming}
%%%%%%%%%%%%%%%%%%%%%%%%%%%%%%%%%%%%%%%%%%%%%%%%%%%%%%%%%%%%

Edit distance measures the minimum number or cost of operations needed to
transform one string into another.

Common operations include:

\[
\boxed{
\text{insertion},
}
\]

\[
\boxed{
\text{deletion},
}
\]

and

\[
\boxed{
\text{substitution}.
}
\]

%%%%%%%%%%%%%%%%%%%%%%%%%%%%%%%%%%%%%%%%%%%%%%%%%%%%%%%%%%%%
\subsection{Edit-Distance Recurrence}
\label{subsec:edit-distance-recurrence}
%%%%%%%%%%%%%%%%%%%%%%%%%%%%%%%%%%%%%%%%%%%%%%%%%%%%%%%%%%%%

Let

\[
D(i,j)
\]

be the edit distance between prefixes of lengths \(i\) and \(j\).

With unit costs,

\[
\boxed{
D(i,j)
=
\min
\begin{cases}
D(i-1,j)+1,\\
D(i,j-1)+1,\\
D(i-1,j-1)+\delta_{ij},
\end{cases}
}
\]

where

\[
\boxed{
\delta_{ij}
=
\begin{cases}
0,
&
x_i=y_j,
\\
1,
&
x_i\neq y_j.
\end{cases}
}
\]

%%%%%%%%%%%%%%%%%%%%%%%%%%%%%%%%%%%%%%%%%%%%%%%%%%%%%%%%%%%%
\subsection{Edit-Distance Boundary Conditions}
\label{subsec:edit-distance-boundary}
%%%%%%%%%%%%%%%%%%%%%%%%%%%%%%%%%%%%%%%%%%%%%%%%%%%%%%%%%%%%

Transforming the first \(i\) symbols into an empty string requires \(i\)
deletions:

\[
\boxed{
D(i,0)=i.
}
\]

Transforming an empty string into the first \(j\) symbols requires \(j\)
insertions:

\[
\boxed{
D(0,j)=j.
}
\]

%%%%%%%%%%%%%%%%%%%%%%%%%%%%%%%%%%%%%%%%%%%%%%%%%%%%%%%%%%%%
\subsection{Sequence Alignment}
\label{subsec:sequence-alignment-dp}
%%%%%%%%%%%%%%%%%%%%%%%%%%%%%%%%%%%%%%%%%%%%%%%%%%%%%%%%%%%%

Sequence alignment generalizes edit distance by assigning application-
specific scores or penalties to:

\[
\boxed{
\text{matches},
}
\]

\[
\boxed{
\text{mismatches},
}
\]

and

\[
\boxed{
\text{gaps}.
}
\]

Dynamic programming is fundamental in computational biology for aligning
DNA, RNA, and protein sequences.

%%%%%%%%%%%%%%%%%%%%%%%%%%%%%%%%%%%%%%%%%%%%%%%%%%%%%%%%%%%%
\subsection{Matrix-Chain Multiplication}
\label{subsec:matrix-chain-multiplication}
%%%%%%%%%%%%%%%%%%%%%%%%%%%%%%%%%%%%%%%%%%%%%%%%%%%%%%%%%%%%

Suppose matrices

\[
A_1A_2\cdots A_n
\]

must be multiplied.

Matrix multiplication is associative, but different parenthesizations
require different numbers of scalar multiplications.

Dynamic programming finds the cheapest parenthesization.

%%%%%%%%%%%%%%%%%%%%%%%%%%%%%%%%%%%%%%%%%%%%%%%%%%%%%%%%%%%%
\subsection{Matrix-Chain State}
\label{subsec:matrix-chain-state}
%%%%%%%%%%%%%%%%%%%%%%%%%%%%%%%%%%%%%%%%%%%%%%%%%%%%%%%%%%%%

Suppose matrix \(A_i\) has dimensions

\[
p_{i-1}\times p_i.
\]

Let

\[
M(i,j)
\]

be the minimum number of scalar multiplications required to compute

\[
A_iA_{i+1}\cdots A_j.
\]

%%%%%%%%%%%%%%%%%%%%%%%%%%%%%%%%%%%%%%%%%%%%%%%%%%%%%%%%%%%%
\subsection{Matrix-Chain Recurrence}
\label{subsec:matrix-chain-recurrence}
%%%%%%%%%%%%%%%%%%%%%%%%%%%%%%%%%%%%%%%%%%%%%%%%%%%%%%%%%%%%

Split the product between \(A_k\) and \(A_{k+1}\).

Then

\[
\boxed{
M(i,j)
=
\min_{i\leq k<j}
\left[
M(i,k)
+
M(k+1,j)
+
p_{i-1}p_kp_j
\right].
}
\]

The boundary condition is

\[
\boxed{
M(i,i)=0.
}
\]

%%%%%%%%%%%%%%%%%%%%%%%%%%%%%%%%%%%%%%%%%%%%%%%%%%%%%%%%%%%%
\subsection{Interval Dynamic Programming}
\label{subsec:interval-dynamic-programming}
%%%%%%%%%%%%%%%%%%%%%%%%%%%%%%%%%%%%%%%%%%%%%%%%%%%%%%%%%%%%

Matrix-chain multiplication is an example of interval dynamic programming.

The state describes an interval

\[
[i,j].
\]

The recurrence chooses a split point

\[
k.
\]

Other interval-DP problems include:

\[
\boxed{
\text{polygon triangulation},
}
\]

and certain

\[
\boxed{
\text{parsing problems}.
}
\]

%%%%%%%%%%%%%%%%%%%%%%%%%%%%%%%%%%%%%%%%%%%%%%%%%%%%%%%%%%%%
\subsection{Optimal Binary Search Tree Preview}
\label{subsec:optimal-binary-search-tree-preview}
%%%%%%%%%%%%%%%%%%%%%%%%%%%%%%%%%%%%%%%%%%%%%%%%%%%%%%%%%%%%

Suppose keys have different search probabilities.

An optimal binary search tree minimizes expected search cost.

The recurrence chooses a root for each interval of keys.

Thus the problem has the generic interval structure

\[
\boxed{
\text{best interval solution}
=
\min_{\text{split}}
[
\text{left}
+
\text{right}
+
\text{root cost}
].
}
\]

%%%%%%%%%%%%%%%%%%%%%%%%%%%%%%%%%%%%%%%%%%%%%%%%%%%%%%%%%%%%
\subsection{Subset Dynamic Programming}
\label{subsec:subset-dynamic-programming}
%%%%%%%%%%%%%%%%%%%%%%%%%%%%%%%%%%%%%%%%%%%%%%%%%%%%%%%%%%%%

Some dynamic programs use a subset itself as part of the state.

For a ground set of \(n\) elements, there are

\[
\boxed{
2^n
}
\]

possible subsets.

Subset DP is therefore exponential, but can still be much faster than
enumerating all permutations.

%%%%%%%%%%%%%%%%%%%%%%%%%%%%%%%%%%%%%%%%%%%%%%%%%%%%%%%%%%%%
\subsection{Held--Karp Dynamic Programming for TSP}
\label{subsec:held-karp-tsp}
%%%%%%%%%%%%%%%%%%%%%%%%%%%%%%%%%%%%%%%%%%%%%%%%%%%%%%%%%%%%

Fix starting vertex \(1\).

Let

\[
V(S,j)
\]

be the minimum cost of a path that:

\begin{itemize}
    \item starts at vertex \(1\);
    \item visits exactly the vertices in \(S\);
    \item ends at vertex \(j\in S\).
\end{itemize}

Then

\[
\boxed{
V(S,j)
=
\min_{i\in S\setminus\{j\}}
\left[
V(S\setminus\{j\},i)
+
c_{ij}
\right].
}
\]

%%%%%%%%%%%%%%%%%%%%%%%%%%%%%%%%%%%%%%%%%%%%%%%%%%%%%%%%%%%%
\subsection{Held--Karp Boundary Condition}
\label{subsec:held-karp-boundary}
%%%%%%%%%%%%%%%%%%%%%%%%%%%%%%%%%%%%%%%%%%%%%%%%%%%%%%%%%%%%

For a singleton destination,

\[
\boxed{
V(\{j\},j)
=
c_{1j}.
}
\]

After all nonstarting vertices have been visited, return to vertex \(1\):

\[
\boxed{
C^*
=
\min_j
\left[
V(V\setminus\{1\},j)
+
c_{j1}
\right].
}
\]

%%%%%%%%%%%%%%%%%%%%%%%%%%%%%%%%%%%%%%%%%%%%%%%%%%%%%%%%%%%%
\subsection{Held--Karp Complexity}
\label{subsec:held-karp-complexity}
%%%%%%%%%%%%%%%%%%%%%%%%%%%%%%%%%%%%%%%%%%%%%%%%%%%%%%%%%%%%

The Held--Karp algorithm uses roughly

\[
\boxed{
O(n2^n)
}
\]

states and examines up to \(n\) predecessors per state.

Its time complexity is therefore approximately

\[
\boxed{
O(n^2 2^n).
}
\]

This is exponential, but far better than enumerating

\[
\boxed{
(n-1)!
}
\]

possible tours.

%%%%%%%%%%%%%%%%%%%%%%%%%%%%%%%%%%%%%%%%%%%%%%%%%%%%%%%%%%%%
\subsection{State Design}
\label{subsec:dynamic-programming-state-design}
%%%%%%%%%%%%%%%%%%%%%%%%%%%%%%%%%%%%%%%%%%%%%%%%%%%%%%%%%%%%

The most important modeling question in dynamic programming is often:

\[
\boxed{
\text{What information must the state remember?}
}
\]

A valid state must preserve enough information for optimal future
decisions.

A poor state may:

\[
\boxed{
\text{omit necessary history},
}
\]

or

\[
\boxed{
\text{store unnecessary detail}.
}
\]

%%%%%%%%%%%%%%%%%%%%%%%%%%%%%%%%%%%%%%%%%%%%%%%%%%%%%%%%%%%%
\subsection{Principle of Sufficient State}
\label{subsec:sufficient-state}
%%%%%%%%%%%%%%%%%%%%%%%%%%%%%%%%%%%%%%%%%%%%%%%%%%%%%%%%%%%%

Two histories can be treated as the same state if they have identical
effects on all future feasible decisions and future objective values.

Thus the state is a compressed summary of relevant past information.

This is closely related to the Markov property.

%%%%%%%%%%%%%%%%%%%%%%%%%%%%%%%%%%%%%%%%%%%%%%%%%%%%%%%%%%%%
\subsection{State-Space Size}
\label{subsec:dynamic-programming-state-space-size}
%%%%%%%%%%%%%%%%%%%%%%%%%%%%%%%%%%%%%%%%%%%%%%%%%%%%%%%%%%%%

Suppose a state contains \(d\) components, each taking \(m\) possible
values.

The number of possible states can be as large as

\[
\boxed{
m^d.
}
\]

Thus even dynamic programming can become computationally difficult in high
dimension.

%%%%%%%%%%%%%%%%%%%%%%%%%%%%%%%%%%%%%%%%%%%%%%%%%%%%%%%%%%%%
\subsection{Curse of Dimensionality}
\label{subsec:curse-of-dimensionality}
%%%%%%%%%%%%%%%%%%%%%%%%%%%%%%%%%%%%%%%%%%%%%%%%%%%%%%%%%%%%

The rapid growth of the state space with dimension is called the curse of
dimensionality.

If each of \(d\) state variables is discretized into \(m\) values, then

\[
\boxed{
\text{number of states}
=
m^d.
}
\]

This exponential dependence on \(d\) is a major limitation of classical
dynamic programming.

%%%%%%%%%%%%%%%%%%%%%%%%%%%%%%%%%%%%%%%%%%%%%%%%%%%%%%%%%%%%
\subsection{State Dominance}
\label{subsec:state-dominance}
%%%%%%%%%%%%%%%%%%%%%%%%%%%%%%%%%%%%%%%%%%%%%%%%%%%%%%%%%%%%

Some states can be discarded if another state is guaranteed to be at
least as good in all relevant future respects.

This is a dominance rule.

Dominance can reduce computation by eliminating states that cannot lead to
an optimal solution.

%%%%%%%%%%%%%%%%%%%%%%%%%%%%%%%%%%%%%%%%%%%%%%%%%%%%%%%%%%%%
\subsection{Finite-Horizon Deterministic Control}
\label{subsec:finite-horizon-deterministic-control}
%%%%%%%%%%%%%%%%%%%%%%%%%%%%%%%%%%%%%%%%%%%%%%%%%%%%%%%%%%%%

Consider

\[
\boxed{
x_{t+1}
=
f_t(x_t,u_t),
}
\]

where

\[
x_t
=
\text{state},
\]

and

\[
u_t
=
\text{control}.
\]

Suppose the total cost is

\[
\boxed{
\sum_{t=0}^{T-1}
\ell_t(x_t,u_t)
+
g(x_T).
}
\]

Then

\[
\boxed{
V_T(x)=g(x),
}
\]

and

\[
\boxed{
V_t(x)
=
\min_u
\left[
\ell_t(x,u)
+
V_{t+1}
(
f_t(x,u)
)
\right].
}
\]

%%%%%%%%%%%%%%%%%%%%%%%%%%%%%%%%%%%%%%%%%%%%%%%%%%%%%%%%%%%%
\subsection{Policy}
\label{subsec:dynamic-programming-policy}
%%%%%%%%%%%%%%%%%%%%%%%%%%%%%%%%%%%%%%%%%%%%%%%%%%%%%%%%%%%%

A policy specifies what action should be taken in each state.

A deterministic policy has form

\[
\boxed{
u_t
=
\pi_t(x_t).
}
\]

An optimal policy satisfies

\[
\boxed{
\pi_t^*(x)
\in
\operatorname*{arg\,min}_u
\left[
\ell_t(x,u)
+
V_{t+1}(f_t(x,u))
\right].
}
\]

%%%%%%%%%%%%%%%%%%%%%%%%%%%%%%%%%%%%%%%%%%%%%%%%%%%%%%%%%%%%
\subsection{Open-Loop Versus Feedback Decisions}
\label{subsec:open-loop-feedback}
%%%%%%%%%%%%%%%%%%%%%%%%%%%%%%%%%%%%%%%%%%%%%%%%%%%%%%%%%%%%

An open-loop plan chooses the full decision sequence in advance.

A feedback policy chooses decisions based on the current state.

Dynamic programming naturally produces feedback policies:

\[
\boxed{
u_t=\pi_t(x_t).
}
\]

This is especially valuable when disturbances or uncertainty can change
the state.

%%%%%%%%%%%%%%%%%%%%%%%%%%%%%%%%%%%%%%%%%%%%%%%%%%%%%%%%%%%%
\subsection{Stochastic Dynamic Programming}
\label{subsec:stochastic-dynamic-programming}
%%%%%%%%%%%%%%%%%%%%%%%%%%%%%%%%%%%%%%%%%%%%%%%%%%%%%%%%%%%%

Suppose the next state depends on random information \(\xi_t\):

\[
\boxed{
x_{t+1}
=
f_t
(
x_t,
u_t,
\xi_t
).
}
\]

Then the Bellman recursion uses expected future value:

\[
\boxed{
V_t(x)
=
\min_u
\left[
\ell_t(x,u)
+
\mathbb{E}
\left[
V_{t+1}
(
f_t(x,u,\xi_t)
)
\right]
\right].
}
\]

%%%%%%%%%%%%%%%%%%%%%%%%%%%%%%%%%%%%%%%%%%%%%%%%%%%%%%%%%%%%
\subsection{Discrete Stochastic Transition Probabilities}
\label{subsec:stochastic-transition-probabilities}
%%%%%%%%%%%%%%%%%%%%%%%%%%%%%%%%%%%%%%%%%%%%%%%%%%%%%%%%%%%%

Suppose

\[
P(s'|s,a)
\]

is the probability of moving from state \(s\) to \(s'\) after action
\(a\).

Then

\[
\boxed{
V_t(s)
=
\min_{a\in A(s)}
\left[
c_t(s,a)
+
\sum_{s'}
P(s'|s,a)
V_{t+1}(s')
\right].
}
\]

This is the basic finite-horizon stochastic Bellman recursion.

%%%%%%%%%%%%%%%%%%%%%%%%%%%%%%%%%%%%%%%%%%%%%%%%%%%%%%%%%%%%
\subsection{Markov Decision Processes}
\label{subsec:markov-decision-processes}
%%%%%%%%%%%%%%%%%%%%%%%%%%%%%%%%%%%%%%%%%%%%%%%%%%%%%%%%%%%%

A Markov decision process, abbreviated MDP, is characterized by:

\[
\boxed{
\text{states},
}
\]

\[
\boxed{
\text{actions},
}
\]

\[
\boxed{
\text{transition probabilities},
}
\]

and

\[
\boxed{
\text{rewards or costs}.
}
\]

A standard notation is

\[
\boxed{
(\mathcal{S},\mathcal{A},P,r).
}
\]

%%%%%%%%%%%%%%%%%%%%%%%%%%%%%%%%%%%%%%%%%%%%%%%%%%%%%%%%%%%%
\subsection{Markov Property}
\label{subsec:mdp-markov-property}
%%%%%%%%%%%%%%%%%%%%%%%%%%%%%%%%%%%%%%%%%%%%%%%%%%%%%%%%%%%%

The Markov property means that conditional on the present state and
action, the distribution of the next state does not depend on the earlier
history.

Symbolically,

\[
\boxed{
P
(
S_{t+1}
\mid
S_t,A_t,S_{t-1},A_{t-1},\ldots
)
=
P
(
S_{t+1}
\mid
S_t,A_t
).
}
\]

This is precisely why the state is sufficient for future decision making.

%%%%%%%%%%%%%%%%%%%%%%%%%%%%%%%%%%%%%%%%%%%%%%%%%%%%%%%%%%%%
\subsection{Expected Reward Formulation}
\label{subsec:mdp-expected-reward}
%%%%%%%%%%%%%%%%%%%%%%%%%%%%%%%%%%%%%%%%%%%%%%%%%%%%%%%%%%%%

For reward maximization,

\[
r(s,a)
\]

denotes immediate expected reward.

A finite-horizon recursion is

\[
\boxed{
V_t(s)
=
\max_{a\in A(s)}
\left[
r_t(s,a)
+
\sum_{s'}
P(s'|s,a)
V_{t+1}(s')
\right].
}
\]

%%%%%%%%%%%%%%%%%%%%%%%%%%%%%%%%%%%%%%%%%%%%%%%%%%%%%%%%%%%%
\subsection{Discounted Infinite-Horizon Objective}
\label{subsec:discounted-infinite-horizon}
%%%%%%%%%%%%%%%%%%%%%%%%%%%%%%%%%%%%%%%%%%%%%%%%%%%%%%%%%%%%

For an infinite horizon, a common objective is

\[
\boxed{
\mathbb{E}
\left[
\sum_{t=0}^{\infty}
\gamma^t
r(S_t,A_t)
\right],
}
\]

where

\[
\boxed{
0\leq\gamma<1
}
\]

is the discount factor.

Future rewards are weighted by powers of \(\gamma\).

%%%%%%%%%%%%%%%%%%%%%%%%%%%%%%%%%%%%%%%%%%%%%%%%%%%%%%%%%%%%
\subsection{Bellman Optimality Equation}
\label{subsec:bellman-optimality-equation}
%%%%%%%%%%%%%%%%%%%%%%%%%%%%%%%%%%%%%%%%%%%%%%%%%%%%%%%%%%%%

For a discounted MDP,

\[
\boxed{
V^*(s)
=
\max_{a}
\left[
r(s,a)
+
\gamma
\sum_{s'}
P(s'|s,a)
V^*(s')
\right].
}
\]

This is the Bellman optimality equation.

It is a fixed-point equation for the optimal value function.

%%%%%%%%%%%%%%%%%%%%%%%%%%%%%%%%%%%%%%%%%%%%%%%%%%%%%%%%%%%%
\subsection{Bellman Operator}
\label{subsec:bellman-operator}
%%%%%%%%%%%%%%%%%%%%%%%%%%%%%%%%%%%%%%%%%%%%%%%%%%%%%%%%%%%%

Define the Bellman optimality operator \(T\) by

\[
\boxed{
(TV)(s)
=
\max_a
\left[
r(s,a)
+
\gamma
\sum_{s'}
P(s'|s,a)
V(s')
\right].
}
\]

Then the optimal value satisfies

\[
\boxed{
V^*=TV^*.
}
\]

%%%%%%%%%%%%%%%%%%%%%%%%%%%%%%%%%%%%%%%%%%%%%%%%%%%%%%%%%%%%
\subsection{Contraction Property}
\label{subsec:bellman-contraction}
%%%%%%%%%%%%%%%%%%%%%%%%%%%%%%%%%%%%%%%%%%%%%%%%%%%%%%%%%%%%

For discounted finite-state MDPs with

\[
0\leq\gamma<1,
\]

the Bellman operator is a contraction in the sup norm:

\[
\boxed{
\|TV-TW\|_\infty
\leq
\gamma
\|V-W\|_\infty.
}
\]

Therefore it has a unique fixed point.

This provides the basis for value iteration.

%%%%%%%%%%%%%%%%%%%%%%%%%%%%%%%%%%%%%%%%%%%%%%%%%%%%%%%%%%%%
\subsection{Value Iteration}
\label{subsec:value-iteration}
%%%%%%%%%%%%%%%%%%%%%%%%%%%%%%%%%%%%%%%%%%%%%%%%%%%%%%%%%%%%

Value iteration repeatedly applies the Bellman operator:

\[
\boxed{
V_{k+1}
=
TV_k.
}
\]

Thus,

\[
\boxed{
V_{k+1}(s)
=
\max_a
\left[
r(s,a)
+
\gamma
\sum_{s'}
P(s'|s,a)
V_k(s')
\right].
}
\]

For finite discounted MDPs, this converges to \(V^*\).

%%%%%%%%%%%%%%%%%%%%%%%%%%%%%%%%%%%%%%%%%%%%%%%%%%%%%%%%%%%%
\subsection{Policy Evaluation}
\label{subsec:policy-evaluation}
%%%%%%%%%%%%%%%%%%%%%%%%%%%%%%%%%%%%%%%%%%%%%%%%%%%%%%%%%%%%

Suppose policy \(\pi\) is fixed.

Its value function satisfies

\[
\boxed{
V^\pi(s)
=
r(s,\pi(s))
+
\gamma
\sum_{s'}
P
(
s'|s,\pi(s)
)
V^\pi(s').
}
\]

This is a linear system when the state space is finite.

%%%%%%%%%%%%%%%%%%%%%%%%%%%%%%%%%%%%%%%%%%%%%%%%%%%%%%%%%%%%
\subsection{Policy Improvement}
\label{subsec:policy-improvement}
%%%%%%%%%%%%%%%%%%%%%%%%%%%%%%%%%%%%%%%%%%%%%%%%%%%%%%%%%%%%

Given \(V^\pi\), define an improved policy by

\[
\boxed{
\pi_{\mathrm{new}}(s)
\in
\operatorname*{arg\,max}_a
\left[
r(s,a)
+
\gamma
\sum_{s'}
P(s'|s,a)
V^\pi(s')
\right].
}
\]

This chooses actions greedily with respect to the current policy value.

%%%%%%%%%%%%%%%%%%%%%%%%%%%%%%%%%%%%%%%%%%%%%%%%%%%%%%%%%%%%
\subsection{Policy Iteration}
\label{subsec:policy-iteration}
%%%%%%%%%%%%%%%%%%%%%%%%%%%%%%%%%%%%%%%%%%%%%%%%%%%%%%%%%%%%

Policy iteration alternates:

\[
\boxed{
\text{policy evaluation}
}
\]

and

\[
\boxed{
\text{policy improvement}.
}
\]

The process continues until the policy no longer changes.

For finite discounted MDPs, policy iteration reaches an optimal policy in
finitely many policy improvements.

%%%%%%%%%%%%%%%%%%%%%%%%%%%%%%%%%%%%%%%%%%%%%%%%%%%%%%%%%%%%
\subsection{Action-Value Function}
\label{subsec:action-value-function}
%%%%%%%%%%%%%%%%%%%%%%%%%%%%%%%%%%%%%%%%%%%%%%%%%%%%%%%%%%%%

The action-value function is

\[
\boxed{
Q^\pi(s,a)
=
\text{expected discounted return when action }a
\text{ is taken in state }s
}
\]

followed by policy \(\pi\).

For the optimal function,

\[
\boxed{
Q^*(s,a)
=
r(s,a)
+
\gamma
\sum_{s'}
P(s'|s,a)
V^*(s').
}
\]

%%%%%%%%%%%%%%%%%%%%%%%%%%%%%%%%%%%%%%%%%%%%%%%%%%%%%%%%%%%%
\subsection{Relation Between \(V^*\) and \(Q^*\)}
\label{subsec:v-q-relation}
%%%%%%%%%%%%%%%%%%%%%%%%%%%%%%%%%%%%%%%%%%%%%%%%%%%%%%%%%%%%

The optimal state value is

\[
\boxed{
V^*(s)
=
\max_a
Q^*(s,a).
}
\]

An optimal policy can therefore choose

\[
\boxed{
\pi^*(s)
\in
\operatorname*{arg\,max}_a
Q^*(s,a).
}
\]

This connection leads directly to reinforcement-learning methods.

%%%%%%%%%%%%%%%%%%%%%%%%%%%%%%%%%%%%%%%%%%%%%%%%%%%%%%%%%%%%
\subsection{Worked Example: One-Step Stochastic Decision}
\label{subsec:worked-stochastic-dp}
%%%%%%%%%%%%%%%%%%%%%%%%%%%%%%%%%%%%%%%%%%%%%%%%%%%%%%%%%%%%

\begin{workedexamplebox}{Expected Future Cost}

Suppose from state \(s\), action \(a_1\) has immediate cost \(2\).

It moves to state \(A\) with probability

\[
0.75
\]

and to state \(B\) with probability

\[
0.25.
\]

Suppose

\[
V_{t+1}(A)=4,
\qquad
V_{t+1}(B)=10.
\]

Then the total expected cost of choosing \(a_1\) is

\[
2
+
0.75(4)
+
0.25(10).
\]

Therefore,

\[
2+3+2.5
=
7.5.
\]

Hence,

\begin{answerbox}
\[
\boxed{
Q_t(s,a_1)=7.5.
}
\]
\end{answerbox}

\end{workedexamplebox}

%%%%%%%%%%%%%%%%%%%%%%%%%%%%%%%%%%%%%%%%%%%%%%%%%%%%%%%%%%%%
\subsection{Inventory Control}
\label{subsec:inventory-dynamic-programming}
%%%%%%%%%%%%%%%%%%%%%%%%%%%%%%%%%%%%%%%%%%%%%%%%%%%%%%%%%%%%

Inventory systems are natural dynamic programs.

Let

\[
x_t
=
\text{inventory before ordering},
\]

and

\[
u_t
=
\text{order quantity}.
\]

If demand is \(D_t\), then a simple transition is

\[
\boxed{
x_{t+1}
=
x_t
+
u_t
-
D_t.
}
\]

Costs may include ordering, holding, and shortage penalties.

%%%%%%%%%%%%%%%%%%%%%%%%%%%%%%%%%%%%%%%%%%%%%%%%%%%%%%%%%%%%
\subsection{Inventory Bellman Equation}
\label{subsec:inventory-bellman}
%%%%%%%%%%%%%%%%%%%%%%%%%%%%%%%%%%%%%%%%%%%%%%%%%%%%%%%%%%%%

For finite horizon,

\[
\boxed{
V_t(x)
=
\min_{u\geq0}
\left[
c_{\mathrm{order}}(u)
+
\mathbb{E}
\left[
c_{\mathrm{hold}}
(
x+u-D_t
)
+
V_{t+1}
(
x+u-D_t
)
\right]
\right].
}
\]

This is developed more deeply in the later Inventory Theory section.

%%%%%%%%%%%%%%%%%%%%%%%%%%%%%%%%%%%%%%%%%%%%%%%%%%%%%%%%%%%%
\subsection{Equipment Replacement}
\label{subsec:equipment-replacement-dp}
%%%%%%%%%%%%%%%%%%%%%%%%%%%%%%%%%%%%%%%%%%%%%%%%%%%%%%%%%%%%

Suppose equipment becomes more expensive to maintain as it ages.

At each period, choose:

\[
\boxed{
\text{keep}
}
\]

or

\[
\boxed{
\text{replace}.
}
\]

The state may be equipment age.

Dynamic programming compares immediate replacement cost with expected
future maintenance cost.

%%%%%%%%%%%%%%%%%%%%%%%%%%%%%%%%%%%%%%%%%%%%%%%%%%%%%%%%%%%%
\subsection{Scheduling Connection}
\label{subsec:scheduling-dp-connection}
%%%%%%%%%%%%%%%%%%%%%%%%%%%%%%%%%%%%%%%%%%%%%%%%%%%%%%%%%%%%

Scheduling problems can sometimes be solved by dynamic programming when a
small state captures:

\[
\boxed{
\text{completed jobs},
}
\]

\[
\boxed{
\text{current time},
}
\]

or

\[
\boxed{
\text{remaining resources}.
}
\]

Subset DP can be useful when the number of jobs is modest.

%%%%%%%%%%%%%%%%%%%%%%%%%%%%%%%%%%%%%%%%%%%%%%%%%%%%%%%%%%%%
\subsection{Optimal Control Connection}
\label{subsec:optimal-control-dp-connection}
%%%%%%%%%%%%%%%%%%%%%%%%%%%%%%%%%%%%%%%%%%%%%%%%%%%%%%%%%%%%

Dynamic programming is one of the principal foundations of optimal
control.

For continuous-state systems, the discrete-time Bellman recursion has a
continuous-time analogue known as the Hamilton--Jacobi--Bellman equation.

That theory is developed later in Section~\ref{sec:optimal-control}.

%%%%%%%%%%%%%%%%%%%%%%%%%%%%%%%%%%%%%%%%%%%%%%%%%%%%%%%%%%%%
\subsection{Continuous-State Dynamic Programming}
\label{subsec:continuous-state-dynamic-programming}
%%%%%%%%%%%%%%%%%%%%%%%%%%%%%%%%%%%%%%%%%%%%%%%%%%%%%%%%%%%%

A state need not be discrete.

Suppose

\[
x_t\in\mathbb{R}^n.
\]

Then

\[
\boxed{
V_t(x)
=
\min_u
\left[
\ell_t(x,u)
+
V_{t+1}(f_t(x,u))
\right].
}
\]

The difficulty is that \(V_t\) is now a function over a continuous state
space.

Exact tabulation may be impossible.

%%%%%%%%%%%%%%%%%%%%%%%%%%%%%%%%%%%%%%%%%%%%%%%%%%%%%%%%%%%%
\subsection{Discretization of Continuous States}
\label{subsec:dp-state-discretization}
%%%%%%%%%%%%%%%%%%%%%%%%%%%%%%%%%%%%%%%%%%%%%%%%%%%%%%%%%%%%

One approach approximates a continuous state space by a finite grid.

For example,

\[
x\in[0,1]
\]

might be replaced by

\[
\boxed{
\left\{
0,
\frac1N,
\frac2N,
\ldots,
1
\right\}.
}
\]

A finer grid improves resolution but greatly increases computational
cost.

%%%%%%%%%%%%%%%%%%%%%%%%%%%%%%%%%%%%%%%%%%%%%%%%%%%%%%%%%%%%
\subsection{Approximate Dynamic Programming}
\label{subsec:approximate-dynamic-programming}
%%%%%%%%%%%%%%%%%%%%%%%%%%%%%%%%%%%%%%%%%%%%%%%%%%%%%%%%%%%%

When exact value functions are too expensive to store, one may approximate

\[
\boxed{
V(s)
\approx
\widehat{V}(s;\theta).
}
\]

The approximation may use:

\[
\boxed{
\text{basis functions},
}
\]

\[
\boxed{
\text{regression},
}
\]

\[
\boxed{
\text{trees},
}
\]

or

\[
\boxed{
\text{neural networks}.
}
\]

This is called approximate dynamic programming.

%%%%%%%%%%%%%%%%%%%%%%%%%%%%%%%%%%%%%%%%%%%%%%%%%%%%%%%%%%%%
\subsection{Value-Function Approximation}
\label{subsec:value-function-approximation}
%%%%%%%%%%%%%%%%%%%%%%%%%%%%%%%%%%%%%%%%%%%%%%%%%%%%%%%%%%%%

A linear approximation may have form

\[
\boxed{
\widehat{V}(s)
=
\theta_1\phi_1(s)
+
\cdots
+
\theta_k\phi_k(s),
}
\]

where

\[
\phi_i(s)
\]

are chosen features.

Instead of storing one value for every state, only the coefficients
\(\theta_i\) are stored.

%%%%%%%%%%%%%%%%%%%%%%%%%%%%%%%%%%%%%%%%%%%%%%%%%%%%%%%%%%%%
\subsection{State Aggregation}
\label{subsec:state-aggregation}
%%%%%%%%%%%%%%%%%%%%%%%%%%%%%%%%%%%%%%%%%%%%%%%%%%%%%%%%%%%%

State aggregation groups several similar states into one approximate
state.

If the original state space is huge,

\[
\mathcal{S},
\]

replace it by a smaller set

\[
\boxed{
\widetilde{\mathcal{S}}.
}
\]

This reduces computation at the cost of approximation error.

%%%%%%%%%%%%%%%%%%%%%%%%%%%%%%%%%%%%%%%%%%%%%%%%%%%%%%%%%%%%
\subsection{Rolling-Horizon Decisions}
\label{subsec:rolling-horizon-dp}
%%%%%%%%%%%%%%%%%%%%%%%%%%%%%%%%%%%%%%%%%%%%%%%%%%%%%%%%%%%%

A long planning horizon may be approximated by repeatedly solving a
shorter-horizon problem.

At each time:

\begin{enumerate}
    \item solve a finite-horizon problem;
    \item implement only the first action;
    \item observe the new state;
    \item solve again.
\end{enumerate}

This is a rolling-horizon strategy.

%%%%%%%%%%%%%%%%%%%%%%%%%%%%%%%%%%%%%%%%%%%%%%%%%%%%%%%%%%%%
\subsection{Dynamic Programming Versus Greedy Methods}
\label{subsec:dp-versus-greedy}
%%%%%%%%%%%%%%%%%%%%%%%%%%%%%%%%%%%%%%%%%%%%%%%%%%%%%%%%%%%%

A greedy algorithm commits immediately to a locally best action.

Dynamic programming evaluates the effect of an action on the optimal
future value.

Thus greedy reasoning uses

\[
\boxed{
\text{immediate benefit},
}
\]

while dynamic programming uses

\[
\boxed{
\text{immediate benefit}
+
\text{optimal future benefit}.
}
\]

%%%%%%%%%%%%%%%%%%%%%%%%%%%%%%%%%%%%%%%%%%%%%%%%%%%%%%%%%%%%
\subsection{Dynamic Programming Versus Branch-and-Bound}
\label{subsec:dp-versus-branch-bound}
%%%%%%%%%%%%%%%%%%%%%%%%%%%%%%%%%%%%%%%%%%%%%%%%%%%%%%%%%%%%

Dynamic programming identifies and merges equivalent subproblems.

Branch-and-bound keeps different subproblems separate but prunes them
using bounds.

Thus:

\[
\boxed{
\text{DP exploits repeated states},
}
\]

while

\[
\boxed{
\text{branch-and-bound exploits bounds}.
}
\]

Some algorithms combine both ideas.

%%%%%%%%%%%%%%%%%%%%%%%%%%%%%%%%%%%%%%%%%%%%%%%%%%%%%%%%%%%%
\subsection{Dynamic Programming Versus Divide and Conquer}
\label{subsec:dp-versus-divide-conquer}
%%%%%%%%%%%%%%%%%%%%%%%%%%%%%%%%%%%%%%%%%%%%%%%%%%%%%%%%%%%%

Divide-and-conquer also decomposes a problem.

The distinction is that divide-and-conquer usually creates independent
subproblems.

Dynamic programming is most useful when subproblems overlap.

Thus:

\[
\boxed{
\text{independent subproblems}
\rightarrow
\text{divide and conquer},
}
\]

\[
\boxed{
\text{overlapping subproblems}
\rightarrow
\text{dynamic programming}.
}
\]

%%%%%%%%%%%%%%%%%%%%%%%%%%%%%%%%%%%%%%%%%%%%%%%%%%%%%%%%%%%%
\subsection{Dynamic Programming Versus Linear Programming}
\label{subsec:dp-versus-lp}
%%%%%%%%%%%%%%%%%%%%%%%%%%%%%%%%%%%%%%%%%%%%%%%%%%%%%%%%%%%%

Dynamic programming and linear programming are different mathematical
frameworks, but important problems can sometimes be formulated using
either.

For example:

\[
\boxed{
\text{shortest path}
}
\]

can be solved by:

\[
\boxed{
\text{graph algorithms},
}
\]

\[
\boxed{
\text{dynamic programming},
}
\]

or

\[
\boxed{
\text{linear programming}.
}
\]

Different formulations reveal different structures.

%%%%%%%%%%%%%%%%%%%%%%%%%%%%%%%%%%%%%%%%%%%%%%%%%%%%%%%%%%%%
\subsection{Computational Complexity of DP}
\label{subsec:dynamic-programming-complexity}
%%%%%%%%%%%%%%%%%%%%%%%%%%%%%%%%%%%%%%%%%%%%%%%%%%%%%%%%%%%%

A common complexity estimate is

\[
\boxed{
\text{number of states}
\times
\text{work per state}.
}
\]

If there are \(N\) states and at most \(A\) actions per state, then a
simple finite-horizon computation may require approximately

\[
\boxed{
O(NA).
}
\]

The challenge is often that \(N\) itself can be enormous.

%%%%%%%%%%%%%%%%%%%%%%%%%%%%%%%%%%%%%%%%%%%%%%%%%%%%%%%%%%%%
\subsection{Sparse State Exploration}
\label{subsec:sparse-state-exploration}
%%%%%%%%%%%%%%%%%%%%%%%%%%%%%%%%%%%%%%%%%%%%%%%%%%%%%%%%%%%%

Memoization can be advantageous if only a small portion of the theoretical
state space is reachable.

Rather than filling every table entry, compute only states reached from
the initial state.

This can save substantial time and memory.

%%%%%%%%%%%%%%%%%%%%%%%%%%%%%%%%%%%%%%%%%%%%%%%%%%%%%%%%%%%%
\subsection{State Compression}
\label{subsec:state-compression}
%%%%%%%%%%%%%%%%%%%%%%%%%%%%%%%%%%%%%%%%%%%%%%%%%%%%%%%%%%%%

Sometimes a state can be represented more compactly.

For example, if only the parity of a count matters, storing the full count
is unnecessary.

Good state compression can reduce both:

\[
\boxed{
\text{memory},
}
\]

and

\[
\boxed{
\text{running time}.
}
\]

%%%%%%%%%%%%%%%%%%%%%%%%%%%%%%%%%%%%%%%%%%%%%%%%%%%%%%%%%%%%
\subsection{Reconstructing Solutions}
\label{subsec:reconstructing-dp-solutions}
%%%%%%%%%%%%%%%%%%%%%%%%%%%%%%%%%%%%%%%%%%%%%%%%%%%%%%%%%%%%

Computing the optimal value alone may not identify the actual solution.

To reconstruct decisions, store:

\[
\boxed{
\text{parent states},
}
\]

\[
\boxed{
\text{chosen actions},
}
\]

or

\[
\boxed{
\text{split locations}.
}
\]

Then trace backward from the final state.

%%%%%%%%%%%%%%%%%%%%%%%%%%%%%%%%%%%%%%%%%%%%%%%%%%%%%%%%%%%%
\subsection{Tie Handling}
\label{subsec:dynamic-programming-ties}
%%%%%%%%%%%%%%%%%%%%%%%%%%%%%%%%%%%%%%%%%%%%%%%%%%%%%%%%%%%%

Several actions can produce the same optimal value.

Then

\[
\operatorname*{arg\,min}
\]

or

\[
\operatorname*{arg\,max}
\]

may contain multiple actions.

A solver may:

\[
\boxed{
\text{return one optimum},
}
\]

or

\[
\boxed{
\text{record all optimal choices}.
}
\]

%%%%%%%%%%%%%%%%%%%%%%%%%%%%%%%%%%%%%%%%%%%%%%%%%%%%%%%%%%%%
\subsection{Numerical Issues}
\label{subsec:dynamic-programming-numerical-issues}
%%%%%%%%%%%%%%%%%%%%%%%%%%%%%%%%%%%%%%%%%%%%%%%%%%%%%%%%%%%%

When costs are floating-point numbers, numerical comparisons can be
sensitive to rounding.

In stochastic problems, repeated expectations can also accumulate
roundoff.

Practical implementations may require:

\[
\boxed{
\text{tolerances},
}
\]

and

\[
\boxed{
\text{stable summation or scaling}.
}
\]

%%%%%%%%%%%%%%%%%%%%%%%%%%%%%%%%%%%%%%%%%%%%%%%%%%%%%%%%%%%%
\subsection{Dynamic Programming Workflow}
\label{subsec:dynamic-programming-workflow}
%%%%%%%%%%%%%%%%%%%%%%%%%%%%%%%%%%%%%%%%%%%%%%%%%%%%%%%%%%%%

A practical workflow is:

\begin{enumerate}
    \item identify the stages or recursive structure;
    \item define the state;
    \item verify that the state contains sufficient information;
    \item define feasible actions;
    \item define immediate cost or reward;
    \item define the transition law;
    \item specify terminal or boundary conditions;
    \item derive the Bellman recurrence;
    \item determine a valid computation order;
    \item estimate the number of states;
    \item estimate work per state;
    \item decide between memoization and tabulation;
    \item store decisions if solution reconstruction is required;
    \item exploit dominance and state compression;
    \item verify the recurrence on small instances;
    \item analyze time and memory complexity.
\end{enumerate}

%%%%%%%%%%%%%%%%%%%%%%%%%%%%%%%%%%%%%%%%%%%%%%%%%%%%%%%%%%%%
\subsection{Choosing a Dynamic Programming Formulation}
\label{subsec:choosing-dp-formulation}
%%%%%%%%%%%%%%%%%%%%%%%%%%%%%%%%%%%%%%%%%%%%%%%%%%%%%%%%%%%%

For sequential decisions, use a stage-state formulation.

For string problems, use:

\[
\boxed{
\text{prefix states}.
}
\]

For interval problems, use:

\[
\boxed{
(i,j)
\text{ states}.
}
\]

For resource constraints, include:

\[
\boxed{
\text{remaining capacity or budget}.
}
\]

For path-dependent discrete problems, consider:

\[
\boxed{
\text{subset states}.
}
\]

For stochastic decisions, include enough information to obtain the Markov
property.

%%%%%%%%%%%%%%%%%%%%%%%%%%%%%%%%%%%%%%%%%%%%%%%%%%%%%%%%%%%%
\subsection{Common Errors}
\label{subsec:dynamic-programming-common-errors}
%%%%%%%%%%%%%%%%%%%%%%%%%%%%%%%%%%%%%%%%%%%%%%%%%%%%%%%%%%%%

\subsubsection{Defining a State That Omits Relevant History}

If two histories produce different future possibilities, they cannot be
merged into the same state.

The Bellman recursion would then be invalid.

%%%%%%%%%%%%%%%%%%%%%%%%%%%%%%%%%%%%%%%%%%%%%%%%%%%%%%%%%%%%
\subsubsection{Storing Too Much History}

A state that stores the entire decision history may be correct but
computationally disastrous.

The goal is a sufficient but compact state.

%%%%%%%%%%%%%%%%%%%%%%%%%%%%%%%%%%%%%%%%%%%%%%%%%%%%%%%%%%%%
\subsubsection{Writing a Recurrence Without Base Cases}

A recursive equation is incomplete without terminal or boundary values.

For example,

\[
V(0,b)=0
\]

is essential in knapsack.

%%%%%%%%%%%%%%%%%%%%%%%%%%%%%%%%%%%%%%%%%%%%%%%%%%%%%%%%%%%%
\subsubsection{Using Future Values Before They Are Available}

Bottom-up tabulation requires a computation order compatible with the
recurrence.

Dependencies must already be solved.

%%%%%%%%%%%%%%%%%%%%%%%%%%%%%%%%%%%%%%%%%%%%%%%%%%%%%%%%%%%%
\subsubsection{Confusing Greedy Choice with Bellman Optimization}

Dynamic programming compares:

\[
\boxed{
\text{current cost}
+
\text{optimal future cost}.
}
\]

Choosing only the smallest immediate cost can be incorrect.

%%%%%%%%%%%%%%%%%%%%%%%%%%%%%%%%%%%%%%%%%%%%%%%%%%%%%%%%%%%%
\subsubsection{Reusing an Item in 0--1 Knapsack}

When using a one-dimensional capacity table, updating capacities in
increasing order can accidentally allow repeated use of the same item.

For \(0\)-\(1\) knapsack, update capacity in decreasing order.

%%%%%%%%%%%%%%%%%%%%%%%%%%%%%%%%%%%%%%%%%%%%%%%%%%%%%%%%%%%%
\subsubsection{Calling \(O(nW)\) Polynomial Without Qualification}

For binary-encoded \(W\),

\[
O(nW)
\]

is pseudo-polynomial.

Formal complexity is measured in input bits.

%%%%%%%%%%%%%%%%%%%%%%%%%%%%%%%%%%%%%%%%%%%%%%%%%%%%%%%%%%%%
\subsubsection{Assuming Dynamic Programming Is Always Efficient}

Dynamic programming can still require exponentially many states.

Subset DP and high-dimensional control are important examples.

%%%%%%%%%%%%%%%%%%%%%%%%%%%%%%%%%%%%%%%%%%%%%%%%%%%%%%%%%%%%
\subsubsection{Ignoring the Curse of Dimensionality}

A continuous or multivariable state space can make direct tabulation
impossible.

Approximation may be necessary.

%%%%%%%%%%%%%%%%%%%%%%%%%%%%%%%%%%%%%%%%%%%%%%%%%%%%%%%%%%%%
\subsubsection{Using Expectation in the Wrong Order}

In stochastic dynamic programming, the Bellman equation optimizes over
actions while accounting for the expectation of future value.

The order depends on when uncertainty is observed.

Information timing must be modeled correctly.

%%%%%%%%%%%%%%%%%%%%%%%%%%%%%%%%%%%%%%%%%%%%%%%%%%%%%%%%%%%%
\subsubsection{Using Future Information in a Policy}

A valid decision at stage \(t\) may only depend on information available
at stage \(t\).

Using future random outcomes creates a nonimplementable policy.

%%%%%%%%%%%%%%%%%%%%%%%%%%%%%%%%%%%%%%%%%%%%%%%%%%%%%%%%%%%%
\subsubsection{Confusing State Value and Action Value}

\[
V(s)
\]

values a state under optimal or specified future behavior.

\[
Q(s,a)
\]

values taking a particular action from that state and then continuing.

They are related but not identical.

%%%%%%%%%%%%%%%%%%%%%%%%%%%%%%%%%%%%%%%%%%%%%%%%%%%%%%%%%%%%
\subsubsection{Forgetting the Discount Factor}

In discounted infinite-horizon models,

\[
0\leq\gamma<1
\]

is part of the Bellman equation.

Omitting it changes the optimization problem.

%%%%%%%%%%%%%%%%%%%%%%%%%%%%%%%%%%%%%%%%%%%%%%%%%%%%%%%%%%%%
\subsubsection{Assuming Approximate DP Gives Exact Optimality}

Function approximation, state aggregation, and sampling reduce
computation but introduce approximation error.

Their solutions need not be exactly optimal.

%%%%%%%%%%%%%%%%%%%%%%%%%%%%%%%%%%%%%%%%%%%%%%%%%%%%%%%%%%%%
\subsection{Applications}
\label{subsec:dynamic-programming-applications}
%%%%%%%%%%%%%%%%%%%%%%%%%%%%%%%%%%%%%%%%%%%%%%%%%%%%%%%%%%%%

\begin{applicationbox}

\textbf{Routing.}
Shortest-path and constrained-path problems can be formulated through
Bellman recursions.

\medskip

\textbf{Inventory.}
Ordering decisions depend on current stock, uncertain demand, and future
costs.

\medskip

\textbf{Scheduling.}
Dynamic programs can optimize sequences of jobs when the state dimension
remains manageable.

\medskip

\textbf{Finance.}
Sequential investment, consumption, and portfolio decisions can be
modeled using state-dependent value functions.

\medskip

\textbf{Control Theory.}
Dynamic programming constructs optimal feedback policies for deterministic
and stochastic dynamical systems.

\medskip

\textbf{Machine Learning.}
Reinforcement learning generalizes dynamic-programming ideas when
transition models or value functions must be learned from data.

\medskip

\textbf{Computational Biology.}
Sequence alignment, edit distance, and biological string comparison use
dynamic-programming tables.

\medskip

\textbf{Operations Research.}
Resource allocation, replacement, inventory, routing, and multistage
decision problems are classical dynamic-programming applications.

\end{applicationbox}

%%%%%%%%%%%%%%%%%%%%%%%%%%%%%%%%%%%%%%%%%%%%%%%%%%%%%%%%%%%%
\subsection{Exercises}
\label{subsec:dynamic-programming-exercises}
%%%%%%%%%%%%%%%%%%%%%%%%%%%%%%%%%%%%%%%%%%%%%%%%%%%%%%%%%%%%

\begin{exercise}[1]
Define dynamic programming.
\end{exercise}

\begin{exercise}[1]
Define optimal substructure.
\end{exercise}

\begin{exercise}[1]
Define overlapping subproblems.
\end{exercise}

\begin{exercise}[1]
Define a state.
\end{exercise}

\begin{exercise}[1]
Define an action.
\end{exercise}

\begin{exercise}[1]
Define a value function.
\end{exercise}

\begin{exercise}[1]
State Bellman's principle of optimality.
\end{exercise}

\begin{exercise}[1]
Distinguish memoization from tabulation.
\end{exercise}

\begin{exercise}[2]
Write a Bellman recursion for a deterministic finite-horizon problem with
transition

\[
s_{t+1}=T(s_t,a_t).
\]
\end{exercise}

\begin{exercise}[2]
Explain why backward induction is natural for finite-horizon dynamic
programming.
\end{exercise}

\begin{exercise}[2]
Write the shortest-path Bellman recursion.
\end{exercise}

\begin{exercise}[2]
Write the \(0\)-\(1\) knapsack recurrence.
\end{exercise}

\begin{exercise}[2]
State the knapsack base cases.
\end{exercise}

\begin{exercise}[2]
Explain why \(0\)-\(1\) knapsack memory can be reduced to \(O(W)\).
\end{exercise}

\begin{exercise}[2]
Explain why capacities must be processed downward in the one-array
\(0\)-\(1\) knapsack implementation.
\end{exercise}

\begin{exercise}[2]
Write the longest common subsequence recurrence.
\end{exercise}

\begin{exercise}[2]
Write the edit-distance recurrence.
\end{exercise}

\begin{exercise}[2]
Write the matrix-chain multiplication recurrence.
\end{exercise}

\begin{exercise}[3]
Solve a four-item \(0\)-\(1\) knapsack problem by constructing the full
dynamic-programming table.
\end{exercise}

\begin{exercise}[3]
Recover the selected items from the knapsack table.
\end{exercise}

\begin{exercise}[3]
Compare full-table and one-dimensional knapsack memory requirements.
\end{exercise}

\begin{exercise}[3]
Solve a resource-allocation problem with three activities and a budget of
five units.
\end{exercise}

\begin{exercise}[3]
Compute the LCS length of

\[
X=\texttt{ABCBDAB}
\]

and

\[
Y=\texttt{BDCABA}.
\]
\end{exercise}

\begin{exercise}[3]
Construct an edit-distance table for two short strings.
\end{exercise}

\begin{exercise}[3]
Use dynamic programming to determine the optimal parenthesization of four
matrices.
\end{exercise}

\begin{exercise}[3]
Derive the time complexity of matrix-chain dynamic programming.
\end{exercise}

\begin{exercise}[3]
Write the Held--Karp recurrence for a four-city TSP.
\end{exercise}

\begin{exercise}[3]
Explain why Held--Karp improves on complete permutation enumeration.
\end{exercise}

\begin{exercise}[3]
Derive the approximate \(O(n^2 2^n)\) time complexity of Held--Karp.
\end{exercise}

\begin{exercise}[3]
Construct a finite-horizon deterministic control problem and derive its
Bellman recursion.
\end{exercise}

\begin{exercise}[3]
Given a small transition table, compute an optimal policy by backward
induction.
\end{exercise}

\begin{exercise}[3]
Write a finite-horizon stochastic Bellman equation using transition
probabilities.
\end{exercise}

\begin{exercise}[3]
Compute one Bellman update for a two-state stochastic problem.
\end{exercise}

\begin{exercise}[3]
Write the discounted Bellman optimality equation for a finite MDP.
\end{exercise}

\begin{exercise}[3]
Perform two iterations of value iteration on a small MDP.
\end{exercise}

\begin{exercise}[3]
Evaluate a fixed policy by solving its Bellman equations.
\end{exercise}

\begin{exercise}[3]
Perform one policy-improvement step.
\end{exercise}

\begin{exercise}[3]
Derive the relationship

\[
V^*(s)=\max_aQ^*(s,a).
\]
\end{exercise}

\begin{exercise}[3]
Formulate a simple inventory-control model as a dynamic program.
\end{exercise}

\begin{exercise}[3]
Formulate an equipment-replacement model.
\end{exercise}

\begin{exercise}[3]
Explain the difference between open-loop and feedback decisions.
\end{exercise}

\begin{exercise}[3]
Give an example of a state that contains unnecessary history and compress
it.
\end{exercise}

\begin{exercise}[3]
Give an example where omitting a state variable invalidates the Bellman
recursion.
\end{exercise}

\begin{exercise}[3]
Explain the curse of dimensionality using a state with six components and
ten possible values per component.
\end{exercise}

\begin{exercise}[3]
Compare dynamic programming and branch-and-bound for a discrete problem.
\end{exercise}

\begin{exercise}[3]
Compare dynamic programming and greedy optimization.
\end{exercise}

\begin{exercise}[3]
Explain why memoization may outperform tabulation when many states are
unreachable.
\end{exercise}

\begin{exercise}[4]
Prove Bellman's principle of optimality for a finite deterministic control
problem.
\end{exercise}

\begin{exercise}[4]
Study shortest paths as linear programming and dynamic programming.
\end{exercise}

\begin{exercise}[4]
Prove correctness of the \(0\)-\(1\) knapsack recurrence.
\end{exercise}

\begin{exercise}[4]
Study the FPTAS for \(0\)-\(1\) knapsack.
\end{exercise}

\begin{exercise}[4]
Study Hirschberg's memory-efficient algorithm for longest common
subsequence.
\end{exercise}

\begin{exercise}[4]
Study sequence alignment with affine gap penalties.
\end{exercise}

\begin{exercise}[4]
Study interval dynamic programming for optimal binary search trees.
\end{exercise}

\begin{exercise}[4]
Study subset dynamic programming for Steiner-tree problems.
\end{exercise}

\begin{exercise}[4]
Derive the Held--Karp algorithm carefully.
\end{exercise}

\begin{exercise}[4]
Study state-dominance rules in resource-constrained dynamic programming.
\end{exercise}

\begin{exercise}[4]
Study stochastic shortest-path dynamic programming.
\end{exercise}

\begin{exercise}[4]
Prove that the discounted Bellman operator is a contraction in the sup
norm.
\end{exercise}

\begin{exercise}[4]
Use the contraction mapping theorem to prove uniqueness of \(V^*\).
\end{exercise}

\begin{exercise}[4]
Prove convergence of value iteration for finite discounted MDPs.
\end{exercise}

\begin{exercise}[4]
Study finite termination of policy iteration for finite discounted MDPs.
\end{exercise}

\begin{exercise}[4]
Study average-cost dynamic programming.
\end{exercise}

\begin{exercise}[4]
Study stochastic shortest-path MDPs without discounting.
\end{exercise}

\begin{exercise}[4]
Derive the Hamilton--Jacobi--Bellman equation as a continuous-time limit.
\end{exercise}

\begin{exercise}[4]
Study approximate value iteration.
\end{exercise}

\begin{exercise}[4]
Investigate state aggregation error.
\end{exercise}

\begin{exercise}[4]
Study fitted value iteration.
\end{exercise}

\begin{exercise}[4]
Investigate approximate policy iteration.
\end{exercise}

\begin{exercise}[4]
Study rollout algorithms.
\end{exercise}

\begin{exercise}[4]
Study the relationship between dynamic programming and reinforcement
learning.
\end{exercise}

\begin{exercise}[4]
Study dynamic programming on tree decompositions.
\end{exercise}

%%%%%%%%%%%%%%%%%%%%%%%%%%%%%%%%%%%%%%%%%%%%%%%%%%%%%%%%%%%%
\subsection{Conceptual Summary}
\label{subsec:dynamic-programming-summary}
%%%%%%%%%%%%%%%%%%%%%%%%%%%%%%%%%%%%%%%%%%%%%%%%%%%%%%%%%%%%

Dynamic programming solves optimization problems by decomposing them into
smaller state-dependent problems.

The key objects are

\[
\boxed{
\text{state}
+
\text{action}
+
\text{transition}
+
\text{value}.
}
\]

For deterministic finite-horizon optimization,

\[
\boxed{
V_t(s)
=
\min_{a\in A_t(s)}
\left[
c_t(s,a)
+
V_{t+1}
(
T_t(s,a)
)
\right].
}
\]

This is the Bellman recursion.

Important discrete examples include:

\[
\boxed{
\text{shortest path},
}
\]

\[
\boxed{
\text{knapsack},
}
\]

\[
\boxed{
\text{sequence alignment},
}
\]

\[
\boxed{
\text{matrix-chain multiplication},
}
\]

and

\[
\boxed{
\text{TSP subset DP}.
}
\]

For stochastic systems,

\[
\boxed{
V_t(s)
=
\min_a
\left[
c_t(s,a)
+
\sum_{s'}
P(s'|s,a)V_{t+1}(s')
\right].
}
\]

For discounted infinite-horizon reward maximization,

\[
\boxed{
V^*(s)
=
\max_a
\left[
r(s,a)
+
\gamma
\sum_{s'}
P(s'|s,a)
V^*(s')
\right].
}
\]

Dynamic programming can dramatically reduce repeated computation, but its
state space can grow exponentially with dimension.

Thus dynamic programming combines

\[
\boxed{
\text{optimization}
+
\text{recursion}
+
\text{algorithms}
+
\text{probability}
+
\text{control}.
}
\]

%%%%%%%%%%%%%%%%%%%%%%%%%%%%%%%%%%%%%%%%%%%%%%%%%%%%%%%%%%%%
\subsection{Section Connections}
\label{subsec:dynamic-programming-connections}
%%%%%%%%%%%%%%%%%%%%%%%%%%%%%%%%%%%%%%%%%%%%%%%%%%%%%%%%%%%%

\chapterconnection{
Dynamic Programming develops recursive optimization through Bellman's
principle of optimality. It connects the discrete search problems of
Section~\ref{sec:combinatorial-optimization} with sequential decision
making, stochastic systems, and feedback control. Knapsack, shortest
paths, sequence problems, and subset optimization illustrate deterministic
dynamic programming, while Markov decision processes introduce stochastic
Bellman equations. These ideas become especially important in later
sections on Stochastic Optimization, Scheduling Theory, Inventory Theory,
Decision Theory, and Optimal Control. The next section, Stochastic
Optimization, develops optimization models in which uncertain quantities
enter objectives, constraints, or future decisions explicitly.
}

%%%%%%%%%%%%%%%%%%%%%%%%%%%%%%%%%%%%%%%%%%%%%%%%%%%%%%%%%%%%
% SECTION SOURCE TRACKING
%%%%%%%%%%%%%%%%%%%%%%%%%%%%%%%%%%%%%%%%%%%%%%%%%%%%%%%%%%%%

%------------------------------------------------------------
% 10.7 Dynamic Programming
%------------------------------------------------------------
%
% Source basis:
%   - Standard dynamic-programming theory.
%   - Standard algorithms and discrete optimization methods.
%   - Standard finite-horizon and Markov decision-process theory.
%
% Used for:
%   - optimal substructure;
%   - overlapping subproblems;
%   - memoization and tabulation;
%   - states, actions, stages, and transitions;
%   - value functions;
%   - Bellman's principle of optimality;
%   - finite-horizon Bellman recursions;
%   - backward induction;
%   - policy reconstruction;
%   - state-space graph interpretation;
%   - shortest-path dynamic programming;
%   - Bellman--Ford connection;
%   - resource allocation;
%   - 0--1 and unbounded knapsack;
%   - pseudo-polynomial complexity;
%   - longest common subsequence;
%   - edit distance;
%   - sequence alignment;
%   - matrix-chain multiplication;
%   - interval dynamic programming;
%   - subset dynamic programming;
%   - Held--Karp TSP dynamic programming;
%   - state design and state compression;
%   - curse of dimensionality;
%   - deterministic control;
%   - policies and feedback;
%   - stochastic dynamic programming;
%   - transition probabilities;
%   - Markov decision processes;
%   - discounted infinite-horizon objectives;
%   - Bellman optimality operators;
%   - contraction property;
%   - value iteration;
%   - policy evaluation and policy iteration;
%   - state-value and action-value functions;
%   - inventory and replacement applications;
%   - continuous-state problems;
%   - state discretization;
%   - approximate dynamic programming;
%   - value-function approximation;
%   - state aggregation;
%   - computational complexity;
%   - applications and exercises.
%
% Notes:
%   - Stochastic Optimization is developed next in Section 10.8.
%   - Inventory Theory is developed later in Section 10.13.
%   - Decision Theory is developed in Section 10.14.
%   - Optimal Control is developed in Section 10.16.
%
%%%%%%%%%%%%%%%%%%%%%%%%%%%%%%%%%%%%%%%%%%%%%%%%%%%%%%%%%%%%